\thispagestyle{empty}

\begin{center}{\normalsize\linespread{1}\selectfont
МИНИСТЕРСТВО НАУКИ И ВЫСШЕГО ОБРАЗОВАНИЯ\\
РОССИЙСКОЙ ФЕДЕРАЦИИ\\
\ \\
федеральное государственное автономное\\
образовательное учреждение высшего высшего образования\\
«Самарский национальный исследовательский университет\\
 имени академика С.П. Королева»\\
 (Самарский университет)\\
\ \\ 
Институт информатики и кибернетики\\
Кафедра технической кибернетики\\
{\fontsize{8pt}{8pt}\selectfont 
\ \\
\ \\
\ \\
\ \\
\ \\
\ \\
\ \\
\ \\
\ \\
\ \\
\ \\
}
ВЫПУСКНАЯ КВАЛИФИКАЦИОННАЯ РАБОТА\\
{\fontsize{8pt}{8pt}\selectfont 
\ \\
\ \\
}
\nohyphens{\fontsize{16pt}{\baselineskip}
\color{red}\bfseries\selectfont 
Исследование погрешности моделирования \\
распространения светового импульса\\
в планарном волноводе}\\
{\fontsize{8pt}{8pt}\selectfont 
\ \\
\ \\
\ \\
}
по программе бакалавриата по направлению подготовки\\
01.03.02  Прикладная математика и информатика,\\
профиль «Компьютерные науки»\\
}


\end{center}
\normalsize 

{\fontsize{8pt}{8pt}\selectfont 
\ \\
\ \\

}
{\linespread{1}\selectfont
\advance\leftskip 2cm
\advance\rightskip 1cm
\noindent
Обучающийся \hrulefill \ \ И.О. Фамилия\\
Научный руководитель ВКР,\\
{\itshape\color{red} должность, ученая степень,\\
ученое звание} \hrulefill \ \ И.О. Фамилия\\
{\color{red}Консультант 
\itshape\fontsize{10pt}{10pt}\selectfont
(\underline{только если назначен})} \hrulefill \ \ И.О. Фамилия\par
}

% Нормоконтролёр вынесен отдельно, поскольку в файле-шаблон у него другой отступ справа
{
\linespread{1}\selectfont
\advance\leftskip 2cm
\advance\rightskip 1.25cm
\noindent
Нормоконтролёр \hrulefill \ С.В. Суханов\par
}

\vfill
\begin{center}
Самара  2023
\end{center}

\clearpage
